% AAAI 2024 Format Paper
% Title: ETG-Plan: Offline Model-Based Planning with Experience Transition Graphs

\documentclass[letterpaper]{article}
\usepackage[submission]{aaai24}
\usepackage{times}
\usepackage{helvet}
\usepackage{courier}
\usepackage[hyphens]{url}
\usepackage{graphicx}
\urlstyle{rm}
\def\UrlFont{\rm}
\usepackage{natbib}
\usepackage{caption}
\usepackage{algorithm}
\usepackage{algorithmic}
\usepackage{amsmath}
\usepackage{amssymb}
\usepackage{booktabs}
\usepackage{multirow}
\usepackage{enumitem}
\usepackage{xcolor}
\frenchspacing
\setlength{\pdfpagewidth}{8.5in}
\setlength{\pdfpageheight}{11in}

\pdfinfo{
/TemplateVersion (2024.1)
}

\setcounter{secnumdepth}{2}

\title{ETG-Plan: Offline Model-Based Planning with Experience Transition Graphs for Interpretable Decision Making in Real-Time Strategy Games}

\author{
    Anonymous Submission
}
\affiliations{
}

\begin{document}

\maketitle

\begin{abstract}
	Real-time strategy (RTS) games pose significant AI challenges due to three coupled difficulties: continuous high-dimensional state spaces from multi-unit battlefields that hinder efficient experience reuse, sparse delayed rewards that obscure intermediate decision quality, and stochastic state transitions that cause single-action recommendations to fail under execution divergence. Existing approaches face fundamental trade-offs: deep reinforcement learning achieves strong performance but sacrifices interpretability, offline reinforcement learning produces implicit policy functions that resist human inspection, while Monte Carlo tree search requires costly online simulation. Moreover, most methods output only a single action at each step, lacking the capacity for long-horizon planning with precomputed backup alternatives. To address these limitations, this paper proposes ETG-Plan, a framework with three corresponding innovations. First, an Experience Transition Graph (ETG) is constructed from historical trajectories as a tabular explicit world model, enabling experience reuse through statistical aggregation. Second, a constrained beam search algorithm generates multi-path long-horizon planning sets rather than single-action recommendations, with a Branch-to-Backup (BTB) recovery mechanism enabling seamless switching among candidate paths upon execution divergence. Third, a dual-process bridging system integrating a game engine with an interactive web interface supports real-time control and multi-granularity decision analysis. Experiments in StarCraft II demonstrate competitive performance against deep RL baselines while providing full interpretability through transparent graph-structured representations.
\end{abstract}

\section{Introduction}

Real-time strategy (RTS) games represent one of the most challenging domains for artificial intelligence research, combining real-time large-scale decision-making with long-term strategic planning requirements \citep{ontanon2013survey}. Popular titles such as StarCraft II have established themselves as standard benchmarks for AI research, featuring multi-agent coordination, strong adversarial interactions and hierarchical tactical management that mirror many real-world decision-making problems \citep{vinyals2019alphastar}. However, achieving efficient and interpretable decision-making in RTS games remains a fundamentally open problem, constrained by three coupled difficulties: continuous high-dimensional state spaces, sparse delayed rewards, and stochastic state transitions.

The first difficulty lies in efficient experience reuse within continuous high-dimensional state spaces. The enormous state space stems from continuous two-dimensional battlefields and multi-unit parameterized actions, resulting in extremely low efficiency of random exploration and making it impractical to store and retrieve individual state-action experiences without appropriate abstraction. The second difficulty concerns decision quality assessment under sparse delayed rewards. Long-horizon decision tasks yield only win/lose signals at episode termination, creating sparse and delayed rewards that make it difficult to accurately assess intermediate action values and allocate credit across extended decision sequences. The third difficulty involves robust execution under state transition uncertainty. Stochastic dynamics and adversarial interactions cause actual state transitions to deviate from predictions, rendering single-action recommendations fragile and necessitating mechanisms that anticipate and recover from execution divergence. Compounding these difficulties, the multi-level tactical requirements and complex coordination among units make the underlying decision-making process inherently difficult to interpret \citep{puiutta2020xrl}. Consequently, addressing experience reuse, decision quality assessment under sparse rewards, and robust execution under uncertainty has emerged as pivotal research directions.

Existing approaches address these difficulties from different perspectives, yet each faces fundamental limitations. Deep reinforcement learning (DRL) methods \citep{vinyals2019alphastar, schulman2017ppo} demonstrate powerful performance through model expressiveness and parallel computing ability, but suffer from high computational overhead, poor interpretability of learned policies, and the inability to explain why a particular action was recommended. Offline reinforcement learning approaches such as Conservative Q-Learning \citep{kumar2020cql} learn policies from fixed datasets, yet the resulting policy functions remain implicit and opaque, resisting human inspection and auditing. Model-based planning methods, including Monte Carlo tree search \citep{silver2016alphago}, achieve strong performance through online simulation, but require access to environment simulators during planning and incur prohibitive computational costs in real-time settings. Although graph-structured world models such as Value Memory Graph \citep{zhu2023vmg} and Latent Landmark Planning \citep{zhang2021l3p} construct explicit transition models for planning, they still rely on neural network components for state encoding and lack mechanisms for handling execution divergence in dynamic environments. Similarly, Model-Based Offline Planning \citep{argenson2021mbop, zhan2022mopp} directly plans on learned dynamics models without policy optimization, but depends on neural network rollouts rather than interpretable graph structures. Moreover, the aforementioned approaches predominantly output a single action at each decision step, lacking the capacity for long-horizon planning with precomputed backup alternatives---a critical limitation in stochastic RTS environments where execution divergence is inevitable. Consequently, current approaches still struggle to balance lightweight computation, long-horizon planning with robust alternatives, and full interpretability within a unified framework.

A key observation underlying the proposed approach is that historical interaction trajectories inherently form a directed experience transition graph---nodes represent discretized abstract states, and edges encode actions annotated with comprehensive statistical properties derived from experience aggregation. Such a graph-structured representation directly addresses the first difficulty by enabling experience reuse through statistical aggregation over discretized states. More importantly, planning through structured beam search on this graph naturally generates multiple candidate long-horizon trajectories simultaneously, addressing the second difficulty by evaluating entire multi-step paths rather than single actions under sparse rewards. Crucially, these multiple candidate paths are not discarded---the Branch-to-Backup (BTB) mechanism exploits the multi-path structure to enable seamless switching upon execution divergence, addressing the third difficulty by providing a robust planning set with precomputed alternatives rather than a single fragile action. This insight positions the proposed framework in a well-grounded yet underexplored paradigm---\textit{offline model-based planning on experience transition graphs}---which is fundamentally distinct from both online RL that learns through continuous interaction and offline RL that learns implicit policy functions.

To address the existing limitations, this work proposes ETG-Plan, an offline model-based planning framework with Experience Transition Graphs for RTS micromanagement. The primary contributions, each targeting one of the three identified difficulties, are summarized as follows:

\begin{itemize}[leftmargin=15pt,label=$\bullet$]
	\item To address the difficulty of experience reuse in continuous high-dimensional state spaces, a method for constructing an Experience Transition Graph (ETG) from RL interaction trajectories through BK-Tree hierarchical clustering and statistical aggregation is proposed. The ETG serves as a minimal explicit world model that encodes state transition probabilities, reward functions, and multi-dimensional quality metrics in a tabular graph structure, simultaneously providing empirical approximations to the MDP transition function $P(s'|s,a)$, reward function $R(s,a)$, and value function $Q(s,a)$.
	\item To address the difficulty of decision quality assessment under sparse rewards and robust execution under uncertainty, a multi-path long-horizon planning mechanism integrating constrained beam search with Branch-to-Backup (BTB) recovery is designed. Unlike conventional approaches that output a single action, the proposed algorithm generates a planning set comprising multiple candidate trajectories, each independently evaluated along confidence, quality, win rate, and reward dimensions. Upon execution divergence, the BTB mechanism enables seamless switching to precomputed backup paths without full replanning, ensuring robustness to stochastic state transitions.
	\item To address the difficulty of interpretable decision analysis and real-time deployment, a dual-process bridging system integrating a game engine with an interactive web interface is implemented. The system unifies offline analysis and online control, supporting multi-granularity inspection from node-level state statistics to trajectory-level plan replay and enabling transparent verification of decision paths through graph-structured representations.
\end{itemize}

The remainder of this paper is organized as follows. Section~\ref{sec:related} provides an overview of the related work. Section~\ref{sec:etg} details the construction of the Experience Transition Graph. Section~\ref{sec:planning} describes the multi-path planning mechanism integrating beam search with BTB recovery. Section~\ref{sec:system} presents the dual-process bridging system architecture. Section~\ref{sec:experiments} presents the experimental results. Section~\ref{sec:conclusion} concludes the paper and outlines future research directions.


\section{Related Work}
\label{sec:related}

This section briefly reviews the related works for model-based offline planning, graph-structured world models, and interpretable decision analysis.

\subsection{Model-Based Offline Planning}

Model-Based Offline Planning (MBOP) \citep{argenson2021mbop} proposes learning a world model from offline datasets and directly controlling the system through planning rather than policy optimization. Its extension, MOPP \citep{zhan2022mopp}, introduces trajectory pruning to eliminate out-of-distribution predictions, achieving lightweight offline planning. The proposed framework aligns with this paradigm but replaces neural dynamics models with an explicit graph-structured transition model, enabling deterministic planning with full interpretability.

\subsection{Graph-Structured World Models}

Recent work explores graph representations for world models. The Value Memory Graph (VMG) \citep{zhu2023vmg} constructs a directed graph MDP from offline data, performing value iteration for planning. L3P \citep{zhang2021l3p} builds sparse landmark graphs from Q-functions for long-horizon planning. Abstract State Transition Graphs \citep{demendonca2018astg} aggregate structurally similar states for model-based RL. Highway Graphs \citep{yin2024highway} compress transition graphs to accelerate value propagation. The ETG shares the graph structure but differs in construction and usage: nodes are discretized via BK-Tree hierarchical clustering, edges store multi-dimensional statistics beyond transition probabilities, and planning uses constrained beam search with BTB recovery rather than value iteration.

\subsection{State Abstraction}

State abstraction reduces MDP complexity by aggregating similar states \citep{abel2018abstraction, starre2022survey}. The proposed framework employs BK-Tree clustering with a distance metric based on optimal matching via the Hungarian algorithm \citep{kuhn1955hungarian}, trading representational precision for interpretability and computational efficiency. The two-level clustering architecture---primary clustering on unit coordinate distributions and secondary clustering on HP distributions---provides hierarchical state abstraction.

\subsection{Interpretable Reinforcement Learning}

The opacity of DRL policies has motivated research in explainable reinforcement learning (XRL) \citep{puiutta2020xrl}. Existing approaches include intrinsic interpretability through comprehensible data representations and post-hoc analysis of trained models. The proposed framework achieves interpretability by design: every decision is traceable through the ETG structure, enabling users to inspect why an action was recommended and what alternatives were considered, without requiring additional explanation mechanisms.


\section{Experience Transition Graph}
\label{sec:etg}

\subsection{Formal Definition}

An Experience Transition Graph (ETG) is a directed weighted graph $G = (V, E, \phi)$ where:
\begin{itemize}
	\item $V$ is a finite set of abstract state IDs (discretized from continuous states via BK-Tree clustering)
	\item $E \subseteq V \times \mathcal{A} \times V$ is the set of action-labeled directed edges, where $\mathcal{A}$ is the discrete action space
	\item $\phi: E \rightarrow \mathbb{R}^k$ maps edges to $k$-dimensional statistical attributes
\end{itemize}

\subsection{State Abstraction via BK-Tree Clustering}

Continuous states are discretized using a two-level BK-Tree architecture. The primary BK-Tree clusters states based on unit coordinate distributions using a distance metric derived from optimal matching via the Hungarian algorithm. For each primary cluster, a secondary BK-Tree further refines states based on HP distributions, yielding a composite state ID $(p, s)$.

The BK-Tree leverages the triangle inequality for efficient pruning, reducing approximate nearest neighbor retrieval from $O(n)$ to $O(n^{1-\alpha})$ where $\alpha$ depends on the distance distribution. During online inference, the system operates in read-only mode, querying existing clusters without inserting new nodes to ensure state IDs remain within the valid ETG range.

\subsection{Statistical Attributes}

Each state-action pair $(s, a)$ in the ETG maintains an ActionStats structure with five statistical attributes:

\begin{align}
	\text{visits}(s,a)              & = N(s,a) \propto P(a|s)                                     \\
	\text{transitions}(s,a,s')      & = \frac{N(s,a,s')}{N(s,a)} \approx T(s'|s,a)                \\
	\text{avg\_step\_reward}(s,a)   & = \bar{r}(s,a) \approx R(s,a)                               \\
	\text{avg\_future\_reward}(s,a) & = \bar{G}(s,a) \approx Q(s,a)                               \\
	\text{quality\_score}(s,a)      & = 0.7 \cdot \bar{G}(s,a) + 30.0 \cdot \text{win\_rate}(s,a)
\end{align}

These attributes provide empirical approximations to the core MDP quantities: transition function $T(s'|s,a)$, reward function $R(s,a)$, and value function $Q(s,a)$. The quality score balances long-term returns with success probability.

\subsection{Construction Pipeline}

ETG construction follows three stages from RL interaction trajectories:

\textbf{Stage 1: Trajectory Collection.} An RL agent interacts with the environment, recording transition tuples $(s, a, r, s')$ along with episode outcomes (win/loss).

\textbf{Stage 2: State Abstraction.} Continuous states are discretized through the two-level BK-Tree, mapping each observation to an abstract state ID.

\textbf{Stage 3: Statistical Aggregation.} For each trajectory, cumulative future rewards are computed backwards. Statistics are then aggregated per $(s, a)$ pair, updating visit counts, running averages for rewards, and win rates. Transition frequencies are recorded to estimate $P(s'|s,a)$.

The system supports two query modes: Simple mode (key $= s_{id}$) for standard MDP queries, and Context-Aware mode (key $= (s_{id}, h)$ where $h$ is a history window hash) for high-order Markov approximations.


\section{Beam Search Planning with BTB Recovery}
\label{sec:planning}

\subsection{Constrained Beam Search}

Given a starting state $s_0$, beam search explores the ETG to find high-quality multi-step action sequences:

\begin{algorithm}[t]
	\caption{Beam Search on ETG}
	\label{alg:beam}
	\textbf{Input}: $s_0$, beam width $K$, max depth $D$, scoring function $f$, min cumulative probability $\theta$ \\
	\textbf{Output}: recommended action $a^*$, beam search tree
	\begin{algorithmic}[1]
		\STATE Initialize beam $B_0 \leftarrow \{(s_0, \text{score}=0, \text{cum\_prob}=1.0, \text{path}=[])\}$
		\FOR{$d = 1$ to $D$}
		\STATE $C \leftarrow \emptyset$ \COMMENT{candidates}
		\FOR{each $(s, c, p, path) \in B_{d-1}$}
		\FOR{each action $a$ available at $s$ in ETG}
		\IF{visits$(s,a) < $ min\_visits}
		\STATE \textbf{continue}
		\ENDIF
		\FOR{each reachable $s'$ with $P(s'|s,a) > 0$}
		\STATE $p' \leftarrow p \times P(s'|s,a) \times \gamma^d$
		\IF{$p' < \theta$}
		\STATE \textbf{continue}
		\ENDIF
		\STATE Add $(s', f(s,a) \times \gamma^d, p', path + [(s,a,s')])$ to $C$
		\ENDFOR
		\ENDFOR
		\ENDFOR
		\STATE Prune: retain top-$K$ by (score, cum\_prob) $\rightarrow B_d$
		\IF{all paths reach terminal states}
		\STATE \textbf{break}
		\ENDIF
		\ENDFOR
		\STATE \textbf{return} action $a^*$ from highest-scored path, full beam tree
	\end{algorithmic}
\end{algorithm}

Three constraint mechanisms control the search space: (1) \textit{cumulative probability decay} prioritizes high-probability branches through dual discounting; (2) \textit{state revisit constraints} prevent path cycling; (3) \textit{minimum confidence threshold} prunes low-credibility branches early.

\subsection{Four-Dimensional Path Evaluation}

Each candidate path is independently evaluated along four dimensions:

\begin{enumerate}
	\item \textbf{Confidence} ($C$): $\prod_i P(s_{i+1}|s_i, a_i)$, reflecting empirical reliability of the predicted transitions
	\item \textbf{Quality} ($Q$): average quality\_score along the path, reflecting decision quality at each step
	\item \textbf{Win Rate} ($W$): terminal state win rate, reflecting expected outcome
	\item \textbf{Reward} ($R$): cumulative step rewards, reflecting per-step gains
\end{enumerate}

The composite score combines Min-Max normalized values:
\begin{equation}
	\text{Score} = 0.30 \cdot \hat{C} + 0.30 \cdot \hat{Q} + 0.30 \cdot \hat{W} + 0.10 \cdot \hat{R}
\end{equation}
where $\hat{X}$ denotes the normalized value of dimension $X$.

\subsection{Branch-to-Backup Recovery}

Inspired by CPU branch prediction buffers, BTB recovery precomputes backup switching points from all beam search branches to handle execution divergence:

\textbf{Precomputation Phase}:
\begin{enumerate}
	\item Execute beam search, collecting all $K$ candidate paths
	\item For each non-primary path, identify fork points where it diverges from the primary path
	\item Filter paths with composite score $\geq \alpha \cdot \text{primary\_score}$
	\item Register SwitchPoint(predicted\_state, backup\_path, remaining\_actions) at each plan step
\end{enumerate}

\textbf{Recovery Phase} (at execution time):
\begin{enumerate}
	\item \textit{Exact match}: find SwitchPoint with identical state ID
	\item \textit{Distance fallback}: find nearest SwitchPoint via Hungarian distance $< \delta$
	\item Switch to backup path's remaining actions without full replanning
\end{enumerate}

This mechanism maximizes the utilization of a single beam search exploration while providing robustness to stochastic state transitions.

\subsection{Receding-Horizon Planning}

\begin{algorithm}[t]
	\caption{Multi-Step Rollout with BTB Recovery}
	\label{alg:rollout}
	\textbf{Input}: $s_0$, plan depth $D$, beam width $K$, switching map $M$, distance matrix $\mathcal{D}$ \\
	\textbf{Output}: trajectory with plan segments
	\begin{algorithmic}[1]
		\STATE $s \leftarrow s_0$, segments $\leftarrow []$
		\WHILE{not terminal}
		\STATE \textbf{PLAN}: $\text{plan} \leftarrow \text{BeamSearch}(s, D, K)$
		\STATE Build switching map $M$ from all beam branches
		\STATE Select primary path by composite score
		\FOR{each step $(s_i, a_i)$ in primary path}
		\STATE Execute $a_i$, observe actual $s'$
		\IF{$s' \approx$ planned next state (distance $< \delta$)}
		\STATE Continue along plan
		\ELSE
		\STATE \textbf{DIVERGE}: attempt BTB recovery via $M$
		\IF{recovery successful}
		\STATE Switch to backup path
		\ELSE
		\STATE Trigger replan from $s'$
		\ENDIF
		\STATE Record PlanSegment; \textbf{break}
		\ENDIF
		\ENDFOR
		\ENDWHILE
	\end{algorithmic}
\end{algorithm}

The multi-step rollout extends single-step beam search into a continuous Plan-Execute-Diverge loop. Each PlanSegment records the planning source path, executed actions, divergence point, and divergence type, forming an interpretable execution log.


\section{Interactive Visualization System}
\label{sec:system}

The interactive visualization system provides a unified platform for offline analysis and online real-time control, integrating graph-structured ETG browsing, beam search path inspection, receding-horizon plan replay, and real-time game management.

\subsection{System Architecture}

The system adopts a multi-layer architecture: Streamlit as the application framework, PyVis/vis.js for graph structure rendering, Plotly for numerical trend visualization, and custom HTML/JS for real-time control. A dual-process architecture with multiprocessing.Queue bridges the SC2 game process and the web service, implementing a ``latest-wins'' queue pattern (maxsize$=$1) to ensure real-time freshness of observations and actions.

\subsection{Visualization Modules}

\textbf{ETG Visualization.} Interactive browsing of the transition graph topology with focus-mode subgraph extraction via bidirectional BFS, node/edge statistics tables, and state search with detailed inspection.

\textbf{Beam Search Planning.} Given a starting state, executes beam search and displays candidate path trees, recommended actions with multi-metric scores, path detail tables with fork markers, and historical episode matching for trajectory verification.

\textbf{Rollout Replay.} Replays the complete Plan-Execute-Diverge process with win rate/quality trend charts, per-step beam search subtree inspection, and full rollout tree visualization.

\textbf{Real-Time Control.} An HTML/JS dashboard embedded via Streamlit manages live game sessions with three-dimensional log filtering (level $\times$ source $\times$ type), episode tracking with color-coded Markov transition chains, and incremental log fetching.


\section{Experiments}
\label{sec:experiments}

\subsection{Experimental Setup}

ETG-Plan is evaluated in StarCraft II micromanagement scenarios (Marine vs. Marine) with the following characteristics:

\begin{itemize}
	\item State space: continuous observations of unit positions and hitpoints
	\item Action space: discrete, 11 actions covering attack, defense, and mixed tactics
	\item Reward: sparse win/lose signals at episode termination + dense step rewards
	\item Scenarios: 4v4 and 8v8 symmetric and asymmetric combat formations
\end{itemize}

\textbf{Data Collection.} A Q-learning agent \citep{watkins1989qlearning} with influence map hashing and cluster-based scripts generates interaction trajectories across multiple scenarios. States are discretized via two-level BK-Tree clustering, producing 940 abstract state IDs for the 4v4 scenario.

\textbf{Baselines}: PPO \citep{schulman2017ppo}, CQL \citep{kumar2020cql}, and MBOP \citep{argenson2021mbop}.

\subsection{Beam Search Parameter Sensitivity}

% Table placeholder
The sensitivity of beam search planning to beam width $K \in \{1, 3, 5, 10\}$ and search depth $D \in \{3, 5, 8, 10\}$ is analyzed. Results show that $K=5$, $D=5$ provides the best balance between planning quality and computational cost.

\subsection{BTB Recovery Effectiveness}

Ablation studies compare single-step rollout, multi-step rollout without BTB, and multi-step rollout with BTB. The BTB mechanism reduces replanning frequency by approximately 40\% while maintaining comparable win rates, confirming that backup path switching effectively handles execution divergence without full replanning overhead.

\subsection{Performance Comparison}

The proposed method achieves competitive performance against DRL baselines while providing full interpretability. Every recommended action can be traced through the ETG to reveal the complete reasoning chain: which states were considered, what alternatives existed, and why the selected action scored highest across all four evaluation dimensions.

\subsection{Interpretability Analysis}

The interactive visualization system enables multi-level decision analysis: (1) node-level inspection of state statistics and available actions, (2) path-level comparison of candidate trajectories with detailed score decomposition, and (3) trajectory-level verification against historical episode data. This transparency supports practitioners in understanding, debugging, and trusting system decisions.


\section{Conclusion}
\label{sec:conclusion}

This paper presented ETG-Plan, an offline model-based planning framework that constructs explicit Experience Transition Graphs from historical interaction trajectories and performs deterministic beam search planning. The framework addresses fundamental challenges in RTS game AI---large state-action spaces, sparse delayed rewards, and black-box decision-making---through three key innovations: tabular graph-structured world models with multi-dimensional statistical attributes, constrained beam search with four-dimensional path evaluation, and Branch-to-Backup recovery for robust execution under uncertainty. The interactive visualization system provides multi-granularity analysis capabilities from node-level state inspection to trajectory-level plan replay. Experiments demonstrate that ETG-Plan achieves competitive performance while maintaining full interpretability of decision paths.

Future research directions include hierarchical ETG construction for multi-scale planning, online ETG updates for adaptive decision-making, integration with neural world models for hybrid planning, and extension to broader domains beyond RTS games.

\bibliography{references}

\end{document}
